\chapternotnumbered{Conclusion}
\label{chap:Conclusion}

The primary goal of this thesis was to investigate the viability of Smoothed Particle Hydrodynamics (SPH) for predicting orographic flows. As outlined in \Cref{chap:FLOWS}, topography-driven flows are responsible for critical atmospheric phenomena, such as internal gravity waves, downslope windstorms, and clear-air turbulence. However, their complex dynamics make them notoriously difficult to model numerically. In \Cref{chap:MODELLING}, we established the mathematical requirements for capturing these flows, highlighting that any successful numerical method must strictly respect the stratified atmospheric background through a \emph{well-balanced} scheme. 

\Cref{chap:SPH} presented the theoretical background and the main novel contributions of this work. After systematically reviewing the classical, Hamiltonian, and Lagrangian derivations of SPH, we proposed a new density-independent formulation that uses potential temperature. Crucially, we then constructed a completely well-balanced discretisation for this framework, greatly upgrading the method's ability to simulate the emergence of mountain waves.

The practical capabilities of these methods were showcased in \Cref{chap:EXPERIMENTS}. The proposed well-balanced formulation outperformed the standard SPH method. It successfully maintained the stationary background and captured the early states of emergence and upward propagation of mountain waves. However, the long-term success of these simulations was ultimately compromised by the onset of \emph{tensile instability}, a known SPH instability.

Ultimately, this thesis demonstrates that SPH possesses potential for atmospheric modelling, but standard formulations are insufficient. Making SPH a reliable tool for atmospheric dynamics will require the development of advanced specialized techniques. Based on our findings, future work must focus on constructing a well-balanced, density-independent $\delta$-SPH formulation to suppress tensile instability without sacrificing the exact hydrostatic balance.

